\input{template.tex}

\ifPDFTeX
  \renewcommand{\rmdefault}{cmr}
\fi

\setlength{\emergencystretch}{2em}

\def\SheetNo{1}

\begin{document}

\sheettitle{Введение. Учимся доказывать.}

\begin{center}\textbf{Решения}\end{center}

\begin{problems}

\item Докажите, что
\[
1\cdot 1!+2\cdot 2!+\ldots+n\cdot n!=(n+1)!-1.
\]
\par\smallskip\noindent\textbf{Решение.} Докажем по индукции по $n$. При $n=1$ имеем $1\cdot1!=2!-1$. Предположим, что равенство верно при некотором $n$. Тогда
\[
\sum_{k=1}^{n+1}k\,k!
=(n+1)!-1+(n+1)(n+1)!
=(n+2)(n+1)!-1
=(n+2)!-1.
\]
Переход доказан.\par

\item Докажите формулу
\[
1^2+2^2+\ldots+n^2=\frac{n(n+1)(2n+1)}{6}.
\]
\par\smallskip\noindent\textbf{Решение.} Проведём индукцию по $n$. При $n=1$ обе части равны $1$. Если формула верна при $n$, то
\[
\begin{aligned}
\sum_{k=1}^{n+1}k^2
&=\frac{n(n+1)(2n+1)}6+(n+1)^2\\
&=\frac{(n+1)(2n^2+7n+6)}6
=\frac{(n+1)(n+2)(2n+3)}6.
\end{aligned}
\]
Это требуемая формула для $n+1$.\par

\item Найдите ошибку в следующем доказательстве утверждения «все лошади одного цвета». База индукции при $n=1$ очевидна. Пусть в любом множестве из $n$ лошадей все лошади одного цвета. Рассмотрим множество из $n+1$ лошадей $h_1,h_2,\ldots,h_{n+1}$. По предположению индукции лошади $h_1,\ldots,h_n$ одного цвета, и лошади $h_2,\ldots,h_{n+1}$ тоже одного цвета. Так как эти два множества пересекаются, все $n+1$ лошадей одного цвета.
\par\smallskip\noindent\textbf{Решение.} Не работает переход от одной лошади к двум. При $n=1$ рассматриваемые множества состоят из $h_1$ и $h_2$ соответственно и не пересекаются. Поэтому нет общей лошади, через цвет которой можно было бы сравнить цвета двух множеств. Пересечение существует лишь при $n\ge2$, но необходимый переход $1\to2$ уже не обоснован.\par

\item Последовательность задана условиями $u_1=2$, $u_2=3$ и $u_n=3u_{n-1}-2u_{n-2}$ при $n>2$. Докажите, что $u_n=2^{n-1}+1$.
\par\smallskip\noindent\textbf{Решение.} Используем индукцию с двумя начальными случаями. При $n=1,2$ имеем $u_1=2^0+1=2$ и $u_2=2^1+1=3$. Пусть формула верна для $n-1$ и $n-2$, где $n\ge3$. Тогда
\[
u_n=3(2^{n-2}+1)-2(2^{n-3}+1)
=2^{n-1}+1.
\]
Следовательно, формула верна для всех $n\ge1$.\par

\item Докажите для всех $n\in\N$:
\begin{parts}
  \item $6\mid(n^3-n)$;
  \item $30\mid(n^5-n)$.
\end{parts}
\par\smallskip\noindent\textbf{Решение.} \par\textbf{а)} Разложим:
\[
n^3-n=(n-1)n(n+1).
\]
Среди трёх последовательных целых чисел есть число, кратное $3$, и хотя бы одно чётное число. Значит, произведение делится на $6$.
\par\textbf{б)} Имеем
\[
n^5-n=(n-1)n(n+1)(n^2+1).
\]
Первые три множителя дают делимость на $6$. Проверим делимость на $5$. Если $n\equiv0,1,-1\pmod5$, один из первых трёх множителей делится на $5$. Если $n\equiv2,-2\pmod5$, то $n^2+1\equiv4+1\equiv0\pmod5$. Таким образом, произведение делится и на $6$, и на $5$, а значит, на $30$.\par

\item Найдите и докажите формулу для суммы
\[
1\cdot 2^1+2\cdot 2^2+\ldots+n\cdot 2^n.
\]
\par\smallskip\noindent\textbf{Решение.} Формула имеет вид
\[
\boxed{\sum_{k=1}^{n}k2^k=(n-1)2^{n+1}+2}.
\]
Докажем её по индукции. При $n=1$ обе части равны $2$. Если формула верна при $n$, то
\[
\begin{aligned}
\sum_{k=1}^{n+1}k2^k
&=(n-1)2^{n+1}+2+(n+1)2^{n+1}\\
&=2n\,2^{n+1}+2
=n2^{n+2}+2.
\end{aligned}
\]
Получили формулу для $n+1$.\par

\item Докажите, что при любом $n\in\N$:
\begin{parts}
  \item
  \[
  1+\frac12+\frac1{2^2}+\ldots+\frac1{2^n}<2;
  \]
  \item
  \[
  \prod_{k=1}^{n}\left(1+\frac1{2^k}\right)<\frac52.
  \]
\end{parts}
\par\smallskip\noindent\textbf{Решение.} \par\textbf{а)} По формуле суммы геометрической прогрессии
\[
1+\frac12+\cdots+\frac1{2^n}
=\frac{1-2^{-(n+1)}}{1-\frac12}
=2-\frac1{2^n}<2.
\]
В частности, эти частичные суммы стремятся к $2$ снизу.
\par\textbf{б)} Обозначим произведение через $P_n$. Усилим утверждение: при $n\ge2$
\[
P_n\le\frac52\left(1-\frac1{2^n}\right).
\]
Для $n=2$ имеем равенство:
\[
P_2=\frac32\cdot\frac54=\frac{15}{8}
=\frac52\left(1-\frac14\right).
\]
Если усиленная оценка верна для $n$, то
\[
\begin{aligned}
P_{n+1}
&=P_n\left(1+\frac1{2^{n+1}}\right)\\
&\le\frac52\left(1-\frac1{2^n}\right)
             \left(1+\frac1{2^{n+1}}\right)\\
&=\frac52\left(1-\frac1{2^{n+1}}-\frac1{2^{2n+1}}\right)
<\frac52\left(1-\frac1{2^{n+1}}\right).
\end{aligned}
\]
По индукции оценка доказана; она влечёт $P_n<5/2$. При $n=1$ непосредственно $P_1=3/2<5/2$.\par

\item Пусть $x\ne0$ и число $x+x^{-1}$ целое. Докажите, что $x^n+x^{-n}$ целое при любом $n\in\N$.
\par\smallskip\noindent\textbf{Решение.} Положим $s=x+x^{-1}\in\mathbb Z$ и $A_n=x^n+x^{-n}$. Имеем $A_0=2$ и $A_1=s$, то есть два начальных значения целые. Раскрывая произведение, получаем
\[
(x+x^{-1})(x^n+x^{-n})
=x^{n+1}+x^{-(n+1)}+x^{n-1}+x^{-(n-1)}.
\]
Следовательно,
\[
A_{n+1}=sA_n-A_{n-1}\qquad(n\ge1).
\]
Если $A_n$ и $A_{n-1}$ целые, то $A_{n+1}$ также целое. Индукция завершает доказательство.\par

\item Последовательность Фибоначчи задана условиями $F_0=0$, $F_1=1$, $F_{n+1}=F_n+F_{n-1}$. Докажите для $m,n\geqslant1$, что
\[
F_{m+n}=F_{m-1}F_n+F_mF_{n+1},
\]
и получите из этого тождество Кассини.
\par\smallskip\noindent\textbf{Решение.} Зафиксируем $m\ge1$ и докажем формулу сложения индукцией по $n$, начиная с $n=0$. При $n=0$ правая часть равна $F_m$, а при $n=1$ --- $F_{m-1}+F_m=F_{m+1}$.
\par Пусть формула верна для $n$ и $n-1$. Тогда
\[
\begin{aligned}
F_{m+n+1}
&=F_{m+n}+F_{m+n-1}\\
&=F_{m-1}(F_n+F_{n-1})+F_m(F_{n+1}+F_n)\\
&=F_{m-1}F_{n+1}+F_mF_{n+2}.
\end{aligned}
\]
Это формула для $n+1$.
\par Теперь получим тождество Кассини. Применяя доказанную формулу к $F_{2n+1}$ двумя способами, имеем
\[
F_{2n+1}=F_n^2+F_{n+1}^2
=F_{n-1}F_{n+1}+F_nF_{n+2}.
\]
Поэтому
\[
F_n^2-F_{n-1}F_{n+1}
=-\bigl(F_{n+1}^2-F_nF_{n+2}\bigr).
\]
Обозначим левую часть через $D_n$. Получили $D_{n+1}=-D_n$. Поскольку $D_1=1$, по индукции $D_n=(-1)^{n-1}$. Итак,
\[
\boxed{F_{n+1}F_{n-1}-F_n^2=(-1)^n}.
\]\par

\item Докажите, что для любого $n\geqslant3$ единицу можно представить в виде суммы $n$ различных дробей с числителем 1 и натуральными знаменателями.
\par\smallskip\noindent\textbf{Решение.} Докажем по индукции по $n\ge3$. База:
\[
1=\frac12+\frac13+\frac16.
\]
Пусть уже имеется представление
\[
1=\frac1{a_1}+\cdots+\frac1{a_n}
\]
с различными натуральными знаменателями. Ни один $a_i$ не равен $1$, иначе остальные положительные слагаемые не могли бы присутствовать. Разделим равенство на $2$ и добавим $1/2$:
\[
1=\frac12+\frac1{2a_1}+\cdots+\frac1{2a_n}.
\]
Знаменатели $2a_i$ попарно различны и не меньше $4$, поэтому не совпадают с новым знаменателем $2$. Получено требуемое представление с $n+1$ слагаемым.\par

\item Назовём натуральное число ровным, если все цифры его десятичной записи одинаковы. Докажите, что любое $n$-значное натуральное число представимо в виде суммы не более чем $n+1$ ровных чисел.
\par\smallskip\noindent\textbf{Решение.} Обозначим $R_n=\underbrace{11\ldots1}_{n\text{ единиц}}$. Сначала докажем усиленное утверждение: любое целое $0\le x<R_n$ представимо суммой не более $n$ ровных чисел. Ноль считаем пустой суммой.
\par При $n=1$ возможно только $x=0$. Пусть утверждение доказано для $n-1$, где $n\ge2$. Поскольку $R_n=10R_{n-1}+1$, из $x<R_n$ следует $x\le10R_{n-1}$. Если $x=10R_{n-1}$, то
\[
x=9R_{n-1}+R_{n-1},
\]
то есть достаточно двух ровных чисел. Иначе делим $x$ на $R_{n-1}$ с остатком:
\[
x=qR_{n-1}+r,\qquad 0\le q\le9,\quad 0\le r<R_{n-1}.
\]
При $q>0$ первое слагаемое --- ровное число, а при $q=0$ его опускаем. По предположению индукции остаток $r$ представим не более чем $n-1$ ровными числами. Всего получаем не более $n$ слагаемых. Усиленное утверждение доказано.
\par Теперь возьмём исходное $n$-значное число $N$ и разделим его на $R_n$:
\[
N=qR_n+r,\qquad 0\le q\le9,\quad 0\le r<R_n.
\]
Число $qR_n$ при $q>0$ ровное, а остаток по доказанному представим не более чем $n$ ровными числами. Следовательно, для $N$ требуется не более $n+1$ ровных чисел.\par

\item На доске записана строка длины 100, состоящая из нулей и единиц. Разрешается либо заменить первую цифру, либо заменить цифру, стоящую сразу после первой единицы. Докажите, что из любой булевой строки можно получить любую другую булевую строку.
\par\smallskip\noindent\textbf{Решение.} Докажем по индукции по $n$, что любую булеву строку длины $n$ можно превратить в любую другую. При $n=1$ это очевидно: при необходимости меняем единственную цифру.
\par Пусть утверждение верно для длины $n-1$. Возьмём произвольную исходную строку длины $n$ и строку, которую хотим получить. По предположению индукции превратим префикс из первых $n-1$ символов исходной строки в $00\ldots01$, не меняя последний символ. Ходы внутри префикса допустимы и в полной строке: если меняется символ после первой единицы префикса, эта единица является первой и во всей строке.
\par Теперь первая единица стоит на позиции $n-1$, поэтому разрешено изменить последний символ. Сделаем его таким, как в требуемой строке, если он ещё не совпадает с нужным.
\par Снова применим предположение индукции к префиксу длины $n-1$ и превратим его в нужный префикс, не трогая последний символ. Получена вся требуемая строка. Индукция завершена; в частности, утверждение верно при $n=100$.\par

\item Выпуклый $n$-угольник разбит непересекающимися диагоналями на треугольники. Докажите, что в разбиении ровно $n-2$ треугольника и $n-3$ диагонали, независимо от способа разбиения.
\par\smallskip\noindent\textbf{Решение.} Проведём индукцию по $n$. При $n=3$ имеется один треугольник и ни одной диагонали, как и требуется.
\par Если $n>3$, выберем одну из диагоналей данного разбиения. Она делит исходный многоугольник на два многоугольника с $p$ и $q$ вершинами, где $p,q<n$ и $p+q=n+2$. В каждом из них остаётся разбиение на треугольники.
\par По предположению индукции число треугольников равно
\[
(p-2)+(q-2)=n-2,
\]
а число проведённых диагоналей исходного многоугольника, с учётом выбранной разделяющей диагонали, равно
\[
(p-3)+(q-3)+1=n-3.
\]
Следовательно, оба количества не зависят от способа разбиения.\par

\item На плоскости проведено несколько прямых общего положения (никакие две из них не параллельны и никакие три не пересекаются в одной точке).
\begin{parts}
  \item На сколько частей они разбивают плоскость?
  \item Может ли среди частей, на которые они разбивают плоскость, оказаться ровно 239 многоугольников?
\end{parts}
\par\smallskip\noindent\textbf{Решение.} \par\textbf{а)} Пусть $R_n$ --- число частей при $n$ прямых; $R_0=1$. Новая, $(n+1)$-я, прямая пересекает прежние $n$ прямых в $n$ различных точках и делится ими на $n+1$ участок. Каждый участок разрезает одну из прежних областей на две. Поэтому
\[
R_{n+1}=R_n+n+1.
\]
Суммируя приращения, получаем
\[
\boxed{R_n=1+\frac{n(n+1)}2}.
\]
\par\textbf{б)} Здесь нужно считать только ограниченные области: именно они являются многоугольниками. При $n\ge2$ неограниченных областей ровно $2n$. Это видно на достаточно большой окружности, содержащей все точки пересечения прямых: прямые пересекают её в $2n$ точках, а полученные $2n$ дуг соответствуют неограниченным областям. Каждой такой области соответствует одна дуга.
\par Следовательно, число многоугольников равно
\[
B_n=R_n-2n=\frac{(n-1)(n-2)}2.
\]
Если бы $B_n=239$, то
\[
(n-1)(n-2)=478,
\qquad n^2-3n-476=0.
\]
Дискриминант равен $1913$, причём $43^2<1913<44^2$, поэтому целого решения нет. При $n=0,1$ многоугольников также нет. Ответ: \textbf{нет}.\par

\item На фестиваль военно-морской песни приглашены хоры из 100 стран. Каждый хор должен исполнить три песни, после чего сразу уехать домой (не выслушивая следующие песни). Ознакомившись с текстами песен, организаторы обнаружили, что каждая песня оскорбительна для одной из участвующих стран. Докажите, что порядок выступлений можно назначить так, чтобы никто не выслушивал более трех оскорбительных для его страны песен.
\par\smallskip\noindent\textbf{Решение.} Докажем утверждение индукцией по числу $m$ стран. При этом разрешим некоторым песням не быть оскорбительными ни для одной из присутствующих стран: такое усиление нужно для удаления страны. При $m=1$ утверждение очевидно, поскольку всего исполняются три песни.
\par Пусть стран $m$. Число песен, оскорбительных для какой-либо из них, не больше $3m$. Поэтому найдётся страна $C$, для которой во всём репертуаре имеется не более трёх оскорбительных песен.
\par Временно удалим хор этой страны. По предположению индукции остальные хоры можно упорядочить нужным образом; песни, оскорбительные только для $C$, при этом считаются неоскорбительными для оставшихся стран.
\par Вернём хор страны $C$ и поставим его последним. Остальные хоры уедут до его выступления и не услышат новых песен. Сам хор $C$ за весь фестиваль услышит не более трёх оскорбительных песен, поскольку во всём репертуаре их для него не больше трёх. Индукционный переход доказан.\par

\item На листе клетчатой бумаги с размером клетки $1\times1$ изображен прямоугольник. Прямоугольник разбит прямыми, параллельными его сторонам, на некоторое количество маленьких прямоугольников. У каждого маленького прямоугольника длины сторон выражаются целыми числами, при этом длина хотя бы одной его стороны чётна. Докажите, что длина хотя бы одной стороны исходного прямоугольника также является чётным числом.
\par\smallskip\noindent\textbf{Решение.} У каждого маленького прямоугольника одна из сторон чётная, а другая целая, поэтому его площадь --- чётное целое число. Площадь большого прямоугольника равна сумме площадей маленьких, следовательно, она также чётная.
\par Стороны большого прямоугольника целые: каждая из них разбита на отрезки, являющиеся сторонами маленьких прямоугольников и имеющие целые длины. Произведение длин двух сторон чётно, значит, хотя бы одна из них чётная.\par

\item В каждой клетке таблицы $10\times10$ стоит одна из цифр. Изначально во всех клетках таблицы стоят нули. У каждой строки и у каждого столбца есть кнопка, которая увеличивает все числа в этой строке или столбце на 1; при этом девятки заменяются на нули. Петя понажимал на какие-то кнопки. Докажите, что таблицу можно вернуть к первоначальному состоянию, нажав на кнопки:
\begin{parts}
  \item не более 180 раз;
  \item не более 90 раз.
\end{parts}
\par\smallskip\noindent\textbf{Решение.} \par\textbf{а)} Пусть Петя нажал кнопки строк $r_1,\ldots,r_{10}$ раз, а кнопки столбцов --- $c_1,\ldots,c_{10}$ раз. В клетке $(i,j)$ теперь стоит остаток $r_i+c_j$ по модулю $10$.
\par Для каждой строки выберем число дополнительных нажатий $\alpha_i\in\{0,\ldots,9\}$, сравнимое с $-r_i$ по модулю $10$, а для каждого столбца --- $\beta_j\in\{0,\ldots,9\}$, сравнимое с $-c_j$. Тогда все клетки станут нулевыми. Кнопок $20$, каждую нажимаем не более $9$ раз, то есть требуется не более $180$ нажатий.
\par\textbf{б)} Рассмотрим десять вариантов восстановления, задаваемых $t=0,\ldots,9$. В варианте $t$ выберем
\[
\alpha_i(t)\equiv-r_i+t\pmod{10},
\qquad
\beta_j(t)\equiv-c_j-t\pmod{10},
\]
где все числа нажатий берутся от $0$ до $9$. Для каждой клетки
\[
r_i+c_j+\alpha_i(t)+\beta_j(t)\equiv0\pmod{10},
\]
поэтому любой вариант восстанавливает таблицу.
\par При переборе $t$ число нажатий каждой конкретной кнопки принимает все значения $0,1,\ldots,9$ по одному разу. Значит, сумма количеств нажатий по всем десяти вариантам равна
\[
20(0+1+\cdots+9)=20\cdot45=900.
\]
Среднее равно $90$, поэтому хотя бы один вариант требует не более $90$ нажатий.\par

\item Взяли несколько одинаковых равносторонних треугольников. Вершины каждого из них пометили цифрами 1, 2 и 3. Затем их сложили в стопку. Могло ли оказаться, что сумма чисел, находящихся в каждом углу, равна 55?
\par\smallskip\noindent\textbf{Решение.} Пусть в стопке $t$ треугольников. Посчитаем сумму всех написанных чисел двумя способами. По треугольникам она равна
\[
(1+2+3)t=6t.
\]
Если бы сумма в каждом из трёх углов равнялась $55$, та же сумма была бы равна $3\cdot55=165$. Но $165$ не делится на $6$. Получили противоречие. Ответ: \textbf{нет}.\par

\item В олимпиаде участвовали $n$ человек. Все решали одинаковые $k$ задач. Будем называть задачу простой, если её решили больше половины человек. Может ли оказаться, что все задачи были простыми, но каждый человек решил меньше половины задач?
\par\smallskip\noindent\textbf{Решение.} Нет. Обозначим через $S$ общее число решённых задач, считая каждую пару «участник --- решённая им задача» отдельно.
\par Если все $k$ задач простые, каждую решили больше $n/2$ участников. Поэтому $S>kn/2$. С другой стороны, если каждый из $n$ участников решил меньше $k/2$ задач, то $S<nk/2$. Эти два неравенства противоречат друг другу.\par

\item В цирке выступает $n$ клоунов. Каждый клоун оделся, использовав в элементах одежды не менее 5 из имеющихся 12 цветов. Известно, что никакие два клоуна не оделись в одинаковый набор цветов, а также никакой цвет не использован более 20 раз. Найдите наибольшее возможное значение $n$.
\par\smallskip\noindent\textbf{Решение.} Посчитаем число пар «клоун --- использованный им цвет». Каждый из $12$ цветов используется не более чем $20$ клоунами, поэтому таких пар не больше $240$. Каждый клоун использует хотя бы $5$ цветов, откуда
\[
5n\le240,\qquad n\le48.
\]
\par Построим пример на $48$ клоунов. Пронумеруем цвета числами $0,\ldots,11$ по кругу. Возьмём четыре набора
\[
\begin{gathered}
A_1=\{0,1,2,3,4\},\qquad A_2=\{0,1,2,3,5\},\\
A_3=\{0,1,2,3,6\},\qquad A_4=\{0,1,2,3,7\}.
\end{gathered}
\]
Для каждого набора возьмём все $12$ циклических сдвигов по модулю $12$.
\par Проверим, что все $48$ наборов различны. В каждом исходном наборе имеется единственный самый большой промежуток между последовательными выбранными цветами по кругу; его длина равна соответственно $8,7,6,5$. Поэтому разные исходные наборы не переводятся друг в друга сдвигом. Каждый набор имеет $12$ различных сдвигов: сдвиг, сохраняющий его, обязан сохранить единственный наибольший промежуток и его концы, а значит, быть нулевым.
\par В $12$ сдвигах любого одного набора каждый цвет встречается ровно $5$ раз: каждый из пяти исходных цветов при переборе сдвигов ровно один раз попадает в заданную позицию. После объединения четырёх семейств каждый цвет используется ровно $20$ раз. Полученные наборы задают костюмы $48$ клоунов и удовлетворяют всем условиям. Ответ:
\[
\boxed{48}.
\]\par

\item В классе 20 школьников. Было устроено несколько экскурсий, в каждой из которых участвовало хотя бы четверо школьников этого класса. Докажите, что найдётся такая экскурсия, что каждый из участвовавших в ней школьников принял участие по меньшей мере в $1/17$ всех экскурсий.
\par\smallskip\noindent\textbf{Решение.} Пусть всего было $m$ экскурсий. Назовём ученика редким, если он участвовал менее чем в $m/17$ экскурсиях.
\par Сначала докажем, что редких учеников не больше $17$. Предположим, что их хотя бы $18$, и выберем $18$ из них. Суммарно они участвовали менее чем в $18m/17$ экскурсиях, считая участия каждого ученика отдельно. Двое оставшихся учеников дают не больше $2m$ участий. Значит, общее число пар «ученик --- посещённая им экскурсия» меньше
\[
\frac{18m}{17}+2m=\frac{52m}{17}<4m.
\]
Но в каждой из $m$ экскурсий участвовало хотя бы четверо, поэтому таких пар не меньше $4m$. Противоречие.
\par Итак, редких учеников не больше $17$, и каждый посетил менее $m/17$ экскурсий. Поэтому суммарное число их участий меньше $m$: при наличии редких учеников оно меньше $17\cdot m/17=m$, а при их отсутствии равно нулю. Число экскурсий, на которых присутствовал хотя бы один редкий ученик, не превосходит этого числа участий и потому также меньше $m$.
\par Следовательно, найдётся экскурсия без редких учеников. Каждый её участник посетил по меньшей мере $m/17$ экскурсий, что и требовалось доказать.\par

\item Докажите, что в любой компании из 10 человек найдутся либо трое попарно знакомых, либо четверо попарно незнакомых.
\par\smallskip\noindent\textbf{Решение.} Сначала докажем лемму: среди любых шести человек найдутся либо трое попарно знакомых, либо трое попарно незнакомых.
\par Выберем одного человека $V$. Среди остальных пяти есть либо трое его знакомых, либо трое незнакомых. В первом случае, если какие-то двое из выбранной тройки знакомы, вместе с $V$ они образуют тройку попарно знакомых; иначе вся выбранная тройка попарно незнакома. Во втором случае, если какие-то двое из выбранной тройки незнакомы, вместе с $V$ они попарно незнакомы; иначе выбранная тройка попарно знакома. Лемма доказана.
\par Теперь рассмотрим компанию из десяти человек и выберем человека $A$. Если у него есть хотя бы четыре знакомых, возьмём любых четырёх. Если среди них есть знакомая пара, она вместе с $A$ даёт тройку попарно знакомых. Если такой пары нет, выбранные четверо попарно незнакомы.
\par Иначе знакомых у $A$ не больше трёх, следовательно, среди остальных девяти у него найдётся хотя бы шесть незнакомых. По лемме среди этих шести есть либо тройка попарно знакомых, либо тройка попарно незнакомых. В первом случае требуемое уже найдено, а во втором добавим $A$ и получим четверых попарно незнакомых.\par

\item В стае обезьян 80 детёнышей. У любых двух из них есть общий дед и, разумеется, у каждого детёныша два деда. Докажите, что найдётся дед, у которого хотя бы 54 внука.
\par\smallskip\noindent\textbf{Решение.} Построим граф: вершины --- деды, а ребро соединяет двух дедов, если у них есть хотя бы один общий внук. Каждый детёныш соответствует ребру между двумя своими дедами; одному ребру могут соответствовать несколько детёнышей. Любые два различных ребра имеют общую вершину: иначе соответствующие им детёныши не имели бы общего деда.
\par Покажем, что такой граф является либо звездой, либо треугольником. Если все рёбра проходят через одну вершину, получаем звезду. Иначе выберем два различных ребра $AB$ и $AC$ с общей вершиной $A$. Поскольку не все рёбра проходят через $A$, существует ребро без вершины $A$. Чтобы пересекаться и с $AB$, и с $AC$, оно должно быть ребром $BC$.
\par Теперь любое ребро обязано пересекаться со всеми тремя рёбрами $AB$, $AC$, $BC$. Ребро с концом вне множества $\{A,B,C\}$ не может этого сделать: его второй конец должен был бы одновременно принадлежать всем трём рёбрам. Значит, других вершин нет, и граф --- треугольник.
\par В случае звезды центральный дед является дедом всех $80$ детёнышей, и утверждение выполнено. В случае треугольника дедов только трое. Каждый детёныш приходится внуком ровно двум из них, поэтому сумма количеств их внуков равна $2\cdot80=160$. По принципу Дирихле у одного из трёх дедов не меньше
\[
\left\lceil\frac{160}{3}\right\rceil=54
\]
внуков.\par

\item Даны 10 двузначных чисел. Докажите, что среди данных чисел можно выбрать два непустых непересекающихся множества так, чтобы суммы чисел в них были равны. Некоторые числа можно вообще не использовать.
\par\smallskip\noindent\textbf{Решение.} Пронумеруем данные числа $a_1,\ldots,a_{10}$ и рассмотрим суммы по всем подмножествам индексов, включая пустое. Подмножеств $2^{10}=1024$, а каждая сумма --- целое число от $0$ до $10\cdot99=990$. Возможных значений суммы всего $991$.
\par По принципу Дирихле найдутся два различных подмножества $A$ и $B$ с одинаковыми суммами. Удалим общие индексы из обоих. Тогда множества $A\setminus B$ и $B\setminus A$ не пересекаются, а суммы по ним по-прежнему равны.
\par Оба множества непусты: если одно пусто, сумма по другому равна нулю, что при положительности всех чисел возможно только для пустого множества. Тогда было бы $A=B$, вопреки выбору. Полученные два набора удовлетворяют условию.\par

\item В ряд выписаны натуральные числа $a_1,a_2,\ldots,a_7$. Докажите, что можно выбрать некоторые из них и поставить перед каждым выбранным числом $+$ или $-$ так, чтобы значение полученного выражения делилось на 101.
\par\smallskip\noindent\textbf{Решение.} Рассмотрим все $2^7=128$ подмножеств множества индексов $\{1,\ldots,7\}$. Каждому подмножеству $A$ сопоставим остаток суммы $\sum_{i\in A}a_i$ по модулю $101$.
\par Остатков всего $101$, поэтому по принципу Дирихле найдутся различные подмножества $A$ и $B$ с одинаковыми остатками. Следовательно,
\[
\sum_{i\in A}a_i-\sum_{i\in B}a_i\equiv0\pmod{101}.
\]
Слагаемые с общими индексами сократятся, и останется выражение
\[
\sum_{i\in A\setminus B}a_i-\sum_{i\in B\setminus A}a_i.
\]
Оно использует некоторые из исходных чисел со знаками $+$ и $-$ и делится на $101$. Выбрано хотя бы одно число, поскольку $A\ne B$.\par

\item Каждый из учеников класса занимается не более чем в двух кружках, причём для любой пары учеников существует кружок, в котором они занимаются вместе. Докажите, что найдётся кружок, в котором занимается не менее $2/3$ всего класса.
\par\smallskip\noindent\textbf{Решение.} Пусть в классе $n\ge2$ учеников. Каждый посещает хотя бы один кружок, иначе у него не было бы общего кружка с другим учеником. Если какой-то ученик посещает только один кружок, то все остальные также посещают этот кружок, и утверждение выполнено.
\par Остаётся случай, когда каждый ученик посещает ровно два кружка. Это та же конструкция, что в задаче 23: кружкам соответствуют деды, а ученикам --- внуки. Построим граф, вершинами которого являются непустые кружки, а ребро соединяет два кружка, если некоторый ученик посещает их оба. Одному ребру могут соответствовать несколько учеников. Любые два ребра имеют общую вершину, поскольку у любых двух учеников есть общий кружок.
\par По доказанному в задаче 23 такой граф является либо звездой, либо треугольником. В случае звезды все ученики посещают кружок, соответствующий её центру. В случае треугольника непустых кружков три, а сумма количеств их участников равна $2n$, поскольку каждый ученик учтён дважды. Значит, в одном из кружков занимается не менее $2n/3$ учеников.\par

\item В классе 28 учеников. На уроке программирования они делятся на три группы. На уроке английского языка они тоже делятся на три группы, но по-другому. И на уроке физкультуры они делятся на три группы каким-то третьим способом. Докажите, что найдутся хотя бы два ученика, которые на всех трёх занятиях находятся друг с другом в одной группе.
\par\smallskip\noindent\textbf{Решение.} По принципу Дирихле в одной из трёх групп на программировании находится хотя бы $\lceil28/3\rceil=10$ учеников. Рассмотрим только их. На английском языке хотя бы $\lceil10/3\rceil=4$ из них попадут в одну группу. Наконец, на физкультуре хотя бы $\lceil4/3\rceil=2$ из этих четырёх учеников окажутся в одной группе. Эта пара находилась вместе на всех трёх занятиях.\par

\item По краю круглого стола равномерно расставлены таблички с фамилиями дипломатов, участвующих в переговорах. После начала переговоров оказалось, что ни один из дипломатов не сидит против своей таблички. Можно ли повернуть стол так, чтобы по крайней мере два дипломата сидели против своих табличек?
\par\smallskip\noindent\textbf{Решение.} Да. Пусть дипломатов $n$. Рассмотрим $n$ положений стола, получающихся поворотами на $0,1,\ldots,n-1$ мест. Каждый дипломат окажется напротив своей таблички ровно в одном из этих положений. Значит, суммарное число совпадений по всем положениям равно $n$.
\par В исходном положении совпадений нет. Если бы в каждом из остальных $n-1$ положений было не больше одного совпадения, их общее число не превосходило бы $n-1$. Это противоречит полученному числу $n$. Следовательно, при некотором повороте совпадений хотя бы два.\par

\item Каждое натуральное число покрасили в один из трёх цветов: красный, синий или зелёный, причём все три цвета встречаются. Может ли оказаться так, что сумма любых двух чисел разных цветов является числом оставшегося цвета?
\par\smallskip\noindent\textbf{Решение.} Нет. Предположим противное. Без ограничения общности число $1$ красное. Пусть $x$ --- наименьшее синее число, а $y$ --- наименьшее зелёное. Поменяв названия синего и зелёного цветов при необходимости, считаем $x<y$.
\par Число $x+1$ зелёное, поскольку $x$ синее, а $1$ красное. Поэтому $x<y\le x+1$, откуда $y=x+1$. Все числа от $1$ до $x-1$ красные.
\par Если $x\ge3$, то число $2$ красное. Тогда $x+2$ зелёное как сумма синего числа $x$ и красного числа $2$. Но то же число $x+2=(x+1)+1$ должно быть синим как сумма зелёного числа $x+1$ и красного числа $1$. Противоречие. Следовательно, $x=2$, а $y=3$.
\par Итак, $1$ красное, $2$ синее, $3$ зелёное. Из равенства $4=1+3$ следует, что $4$ синее. Теперь $5=2+3$ должно быть красным, а $5=1+4$ --- зелёным. Снова противоречие, значит, требуемой раскраски не существует.\par

\item На столе лежат 15 салфеток произвольной формы и размера, которые его полностью закрывают. Докажите, что можно убрать 8 из них так, чтобы оставшиеся закрывали не менее $7/15$ площади стола.
\par\smallskip\noindent\textbf{Решение.} Примем площадь стола за $1$. Упорядочим салфетки, условно расположив их друг над другом. Каждую точку стола закрепим за самой верхней из покрывающих её салфеток. Так каждой салфетке соответствует её собственный участок стола, не закрытый ни одной из расположенных выше салфеток. Наличие салфеток под этим участком не имеет значения.
\par Полученные участки попарно не пересекаются и вместе составляют весь стол. Обозначим их площади через $s_1,\ldots,s_{15}$ и перенумеруем салфетки так, чтобы
\[
s_1\ge s_2\ge\cdots\ge s_{15}\ge0,
\qquad s_1+\cdots+s_{15}=1.
\]
Покажем, что $s_1+\cdots+s_7\ge7/15$. Предположим обратное. Тогда
\[
s_8+\cdots+s_{15}>\frac8{15},
\]
поэтому хотя бы одно из этих восьми чисел больше $1/15$. Но тогда каждое из первых семи чисел также больше $1/15$, откуда $s_1+\cdots+s_7>7/15$. Противоречие.
\par Оставим салфетки, которым соответствуют первые семь участков, а остальные восемь уберём. Каждый из этих семи участков по-прежнему покрыт своей салфеткой, поэтому оставшиеся салфетки покрывают не менее $s_1+\cdots+s_7\ge7/15$ площади стола.\par

\end{problems}

\end{document}
